My mechanical training and ongoing computer-science degree let me reason across robot geometry, learning, and software, but my experience remains centered on bounded research and laboratory integrations. Hermite taught me to test action representations under a defined offline protocol; CR5 work taught me that synchronized observations and carefully timed execution determine whether a learned command is meaningful on hardware. I now need deployment discipline that can carry such reasoning beyond one laboratory setup and toward the varied constraints of industrial environments. HKUST's MEng in Robotics and Embodied AI addresses that transition through its joint ECE/MAE offering.

For the first cohort planned for Fall 2027, the announced sequence of one academic year of coursework at HKUST followed by a one-year industry internship in year two would form a bridge from technical integration to production readiness. The coursework year could help me connect mechanical constraints, learning decisions, and system interfaces; the internship could then expose those choices to different operational requirements. I would enter that progression able to trace failures across timestamps, action semantics, workspace checks, and controller timing, while remaining ready to learn the validation practices required outside a research prototype. My goal is to become a robotics and embodied-AI R\&D engineer who can move laboratory methods toward dependable real-world systems without confusing prototype evidence with production deployment.
